\section*{Capítulo \ref{ch:b2:plant-life-cycles} --- Ciclos de vida y reproducción de las plantas terrestres}

\begin{solution}{exo:b2:plant-life-cycles:1}
\omterm{def:b2:plant-life-cycles:cycles}{Gametofito}: la generación \omterm{def:b2:unicellular-diversity:unicellular}{pluricelular} \omterm{def:b2:meiosis-heredity:meiosis}{haploide}, formada por mitosis a
partir de una espora, que produce los gametos por mitosis. \omterm{def:b2:plant-life-cycles:cycles}{Esporofito}: la
generación \omterm{def:b2:meiosis-heredity:meiosis}{diploide}, formada por mitosis a partir de un cigoto, que
produce las esporas por \omterm{def:b2:meiosis-heredity:meiosis}{meiosis}. Espora: una célula \omterm{def:b2:meiosis-heredity:meiosis}{haploide} formada por
\omterm{def:b2:meiosis-heredity:meiosis}{meiosis} que se desarrolla sin fusionarse. Gameto: una célula \omterm{def:b2:meiosis-heredity:meiosis}{haploide}
formada por mitosis (en las plantas) que debe fusionarse con otra.
\end{solution}

\begin{solution}{exo:b2:plant-life-cycles:2}
Las células de la hoja son $2n = 1260$: espora 630, célula del \omterm{def:b2:plant-life-cycles:fern}{prótalo}
630, anterozoide 630, oosfera 630, cigoto 1260.
\end{solution}

\begin{solution}{exo:b2:plant-life-cycles:3}
Cojín verde: \omterm{def:b2:plant-life-cycles:cycles}{gametofito}. Pedicelo y cápsula: \omterm{def:b2:plant-life-cycles:cycles}{esporofito}. Fronda de
\omterm{def:b2:plant-life-cycles:fern}{helecho}: \omterm{def:b2:plant-life-cycles:cycles}{esporofito}. \omterm{def:b2:plant-life-cycles:fern}{Prótalo}: \omterm{def:b2:plant-life-cycles:cycles}{gametofito}. Grano de \omterm{def:b2:plant-life-cycles:heterospory}{polen}: \omterm{def:b2:plant-life-cycles:cycles}{gametofito}
masculino. Reserva de alimento de la semilla (en un pino): \omterm{def:b2:plant-life-cycles:cycles}{gametofito}
femenino. Embrión de la semilla: el siguiente \omterm{def:b2:plant-life-cycles:cycles}{esporofito}.
\end{solution}

\begin{solution}{exo:b2:plant-life-cycles:4}
Haplonte (\emph{Chlamydomonas}): \omterm{def:b2:meiosis-heredity:meiosis}{meiosis} en el cigoto, de inmediato.
Diplonte (animales, \emph{Fucus}): la \omterm{def:b2:meiosis-heredity:meiosis}{meiosis} produce los gametos.
Haplodiplonte (\omterm{def:b2:plant-life-cycles:moss}{musgos}, \omterm{def:b2:plant-life-cycles:fern}{helechos}, plantas con semilla, muchas algas): la
\omterm{def:b2:meiosis-heredity:meiosis}{meiosis} produce esporas en los \omterm{def:b2:plant-life-cycles:fern}{esporangios} del \omterm{def:b2:plant-life-cycles:cycles}{esporofito}.
\end{solution}

\begin{solution}{exo:b2:plant-life-cycles:5}
$0.03/10^{-4} = \qty{300}{s}$, cinco minutos. En \qty{20}{min},
\qty{12}{cm}. La fecundación solo funciona entre plantas separadas por
unos pocos centímetros, de modo que los \omterm{def:b2:plant-life-cycles:moss}{musgos} crecen en cojines densos
con ambos sexos (o ambos órganos) al alcance.
\end{solution}

\begin{solution}{exo:b2:plant-life-cycles:6}
$v_s = 2\times(1.5\times 10^{-5})^2\times 1099\times 9.81/(9\times
1.8\times 10^{-5}) = \qty{0.030}{m/s}$, \qty{3}{cm/s}; distancia $uh/v_s
= 2\times 0.5/0.03 = \qty{33}{m}$. Semilla: $2\times 20/0.8 =
\qty{50}{m}$ --- una unidad más pesada, liberada desde más alto y con
un ala, llega más o menos igual de lejos.
\end{solution}

\begin{solution}{exo:b2:plant-life-cycles:7}
$200\times 40\times 64 = \num{512000}$ esporas por fronda; $5.1$
\omterm{def:b2:plant-life-cycles:fern}{prótalos}; $0.1$ \omterm{def:b2:plant-life-cycles:cycles}{esporofitos} jóvenes por fronda.
\end{solution}

\begin{solution}{exo:b2:plant-life-cycles:8}
Una semilla es un \omterm{def:b2:plant-life-cycles:heterospory}{óvulo} retenido: el \omterm{def:b2:plant-life-cycles:cycles}{gametofito} femenino debe quedar
sobre la planta madre, encerrado y alimentado, mientras el \omterm{def:b2:plant-life-cycles:cycles}{gametofito}
masculino debe viajar hasta él. Eso exige dos tipos de espora con dos
destinos. La única espora de una planta homospórica debe liberarse para
dar un \omterm{def:b2:plant-life-cycles:cycles}{gametofito} bisexual libre; retenerla retendría también la función
masculina y obligaría a la autofecundación en todas las plantas, y no
habría nada que viajara.
\end{solution}

\begin{solution}{exo:b2:plant-life-cycles:9}
Espora: $\tfrac43\pi(1.5\times 10^{-5})^3\times 1100 =
\qty{1.6e-11}{kg}$, \qty{16}{ng}; energía $0.5\times 1.6\times
10^{-8}\times 2\times 10^{4} = \qty{1.6e-4}{J}$. Semilla: \qty{5}{mg},
$0.5\times 5\times 10^{-3}\times 2\times 10^{4} = \qty{50}{J}$ ---
$3\times 10^{5}$ veces más. La espora puede formar una pequeña lámina
fotosintética y nada más antes de recibir luz; la semilla puede formar
una raíz y un tallo en la oscuridad, esperar una estación entera y
sobrevivir a ser comida o desecada.
\end{solution}

\begin{solution}{exo:b2:plant-life-cycles:10}
$10^{11}/(10^{6}\times 20) = \num{5000}$ granos por metro cúbico; flujo
$5000\times 10^{-7}\times 2 = \qty{1e-3}{}$ granos por segundo; unos
\qty{1000}{s}, un cuarto de hora, por grano. El cono debe estar abierto y
receptivo durante horas o días, y justo cuando los conos masculinos
liberan el \omterm{def:b2:plant-life-cycles:heterospory}{polen} --- la polinización se sincroniza a la semana.
\end{solution}

\begin{solution}{exo:b2:plant-life-cycles:11}
La diploidía enmascara las \omterm{def:b2:genome-diversification:mutation}{mutaciones} recesivas perjudiciales; un cuerpo
\omterm{def:b2:meiosis-heredity:meiosis}{diploide} con tejido vascular puede ser grande y longevo, y la altura
favorece a la vez la captación de luz y la \omterm{thm:b2:plant-life-cycles:dispersal}{dispersión de las esporas}; las
esporas formadas en lo alto de un \omterm{def:b2:plant-life-cycles:cycles}{esporofito} se dispersan por el aire,
mientras que los gametos tienen que nadar. La generación \omterm{def:b2:meiosis-heredity:meiosis}{haploide}
persiste porque la \omterm{def:b2:meiosis-heredity:meiosis}{meiosis} y la fecundación exigen una fase \omterm{def:b2:meiosis-heredity:meiosis}{haploide}, y
el grano de \omterm{def:b2:plant-life-cycles:heterospory}{polen} y el saco embrionario son los cuerpos mínimos que
llevan los gametos el uno hacia el otro y alimentan al embrión.
\end{solution}

\begin{solution}{exo:b2:plant-life-cycles:12}
Cultivando esporas y siguiéndolas al microscopio pudo establecer que una
espora da un pequeño cuerpo portador de órganos sexuales, que el embrión
nace de la oosfera fecundada dentro del \omterm{def:b2:plant-life-cycles:moss}{arquegonio}, que de él crece la
planta portadora de esporas, y que los mismos órganos existen, reducidos,
en el \omterm{def:b2:plant-life-cycles:heterospory}{óvulo} de las coníferas: la alternancia como sucesión de cuerpos y
de órganos. No podía saber qué distinguía las dos generaciones a escala
celular. Los recuentos de cromosomas mostraron que los dos cuerpos
difieren en un factor dos en su número de cromosomas, que la \omterm{def:b2:meiosis-heredity:meiosis}{meiosis} se
sitúa en la formación de las esporas y la fecundación en la oosfera, y
explicaron así por qué la alternancia es obligada.
\end{solution}

\begin{solution}{pb:b2:plant-life-cycles:1}
\textbf{1.} $10\times 200\times 40\times 64 = 5.1\times 10^{6}$
esporas.
\textbf{2.} $\tfrac43\pi(1.5\times 10^{-5})^3\times 1100 =
\qty{1.6e-11}{kg}$ por espora; cosecha de $8\times 10^{-5}$ kg, \qty{80}{mg}.
\textbf{3.} $v_s = 2\times 2.25\times 10^{-10}\times 1099\times
9.81/(1.62\times 10^{-4}) = \qty{0.030}{m/s}$.
\textbf{4.} $x = 1.5\times 0.5/0.030 = \qty{25}{m}$; disco de
$\pi\times 25^2 = \qty{2000}{m^2}$; unas \num{2600} esporas por metro
cuadrado.
\textbf{5.} Las corrientes ascendentes de unos pocos centímetros por
segundo superan $v_s$, de modo que una espora atrapada en la turbulencia
permanece en el aire durante horas y recorre cientos de kilómetros; una
sola espora funda una población si su \omterm{def:b2:plant-life-cycles:fern}{prótalo} puede autofecundarse (lleva
ambos órganos), y por eso las islas remotas tienen floras de \omterm{def:b2:plant-life-cycles:fern}{helechos}
ricas.
\textbf{6.} Espora: $0.5\times 1.6\times 10^{-8}\times 2\times 10^{4} =
\qty{1.6e-4}{J}$; cosecha: \qty{820}{J}; una semilla de pino: \qty{50}{J}
--- toda la cosecha de esporas equivale a dieciséis semillas.
\textbf{7.} $5.1\times 10^{6}/2\times 10^{4} = 256$ \omterm{def:b2:plant-life-cycles:fern}{prótalos}.
\textbf{8.} $1 - 0.8^{6} = 1 - 0.26 = 0.74$.
\textbf{9.} $256\times 0.74\times 0.5\times 0.02 = 1.9$ \omterm{def:b2:plant-life-cycles:fern}{helechos} adultos
por temporada.
\textbf{10.} $30\times 1.9 \approx 57$: muchos más que el único reemplazo
de una población estable, de modo que las cifras de supervivencia son
optimistas --- en un hábitat saturado la mayoría de los \omterm{def:b2:plant-life-cycles:cycles}{esporofitos}
jóvenes mueren por hacinamiento y sombra.
\textbf{11.} $10^{-3}/10^{-4} = \qty{10}{s}$. Muchos \omterm{def:b2:plant-life-cycles:fern}{helechos} maduran los
\omterm{def:b2:plant-life-cycles:moss}{anteridios} y los \omterm{def:b2:plant-life-cycles:moss}{arquegonios} en momentos distintos, o liberan una hormona
(el anteridiógeno) que vuelve masculinos a los \omterm{def:b2:plant-life-cycles:fern}{prótalos} jóvenes vecinos,
de modo que los anterozoides proceden de otro individuo.
\textbf{12.} Tiene una sola célula de grosor, carece de raíces y de
protección, necesita agua líquida para la fecundación y vive unas pocas
semanas: casi toda la mortalidad del ciclo recae sobre él (una espora de
cada $2\times 10^{4}$ llega a formarlo; una cuarta parte de ellos nunca
tiene una noche húmeda).
\textbf{13.} $v_s = 2\times 6.25\times 10^{-10}\times 599\times
9.81/(1.62\times 10^{-4}) = \qty{0.045}{m/s}$.
\textbf{14.} $3\times 15/0.045 = \qty{1000}{m}$.
\textbf{15.} $10^{11}/(2\times 10^{7}) = \num{5000}$ por metro cúbico.
\textbf{16.} $5000\times 10^{-7}\times 2 = \qty{1e-3}{s^{-1}}$: 3,6 por
hora; el primero llega al cabo de unos \qty{1000}{s}, un cuarto de hora.
\textbf{17.} Cinco kilómetros viento abajo, la nube se ha extendido
lateralmente y hacia arriba y la mayoría de los granos se ha depositado
(un alcance de \qty{1}{km} por cada \qty{15}{m} de altura), de modo que la
concentración es varios órdenes de magnitud menor y un \omterm{def:b2:plant-life-cycles:heterospory}{óvulo} puede
esperar días a recibir un grano; el propio \omterm{def:b2:plant-life-cycles:heterospory}{polen} del árbol da sobre todo
semillas por autofecundación, que abortan. La polinización por el viento
necesita una nube densa; de ahí los rodales.
\textbf{18.} $10^{14}$ granos para $2\times 10^{6}$ \omterm{def:b2:plant-life-cycles:heterospory}{óvulos}: $5\times
10^{7}$ granos por \omterm{def:b2:plant-life-cycles:heterospory}{óvulo}.
\textbf{19.} $2\times 10^{6}$ \omterm{def:b2:plant-life-cycles:heterospory}{óvulos} $\times 0.6\times 0.8 = 9.6\times
10^{5}$ semillas; \qty{4.8}{kg}; \qty{48}{MJ}.
\textbf{20.} $3\times 20/0.8 = \qty{75}{m}$: el \omterm{def:b2:plant-life-cycles:heterospory}{polen} llega a un
kilómetro, la semilla a unas pocas alturas de árbol; los genes viajan con
el \omterm{def:b2:plant-life-cycles:heterospory}{polen}, la planta con la semilla.
\textbf{21.} $9.6\times 10^{5}\times 0.05\times 0.01 = 480$ descendientes
adultos, frente a los 57 del \omterm{def:b2:plant-life-cycles:fern}{helecho}.
\textbf{22.} Pino: $4.8\times 10^{7}/480 = \qty{100}{kJ}$ por descendiente
adulto; \omterm{def:b2:plant-life-cycles:fern}{helecho}: $820/1.9 = \qty{430}{J}$ por descendiente adulto,
doscientas veces menos.
\textbf{23.} Ambos tienen la mitad de materia seca a \qty{20}{kJ/g}, de
modo que, con un metabolismo basal de, por ejemplo, \qty{0.5}{mW} por
gramo seco, ambos duran $10^{4}\,
\text{J/g}/(5\times 10^{-4}\,\text{W/g}) = 2\times 10^{7}$ s, unos
ocho meses: la duración de la latencia no es lo que compra la energía de
la semilla. Lo que compra es lo que ocurre tras la germinación --- una
raíz y un tallo construidos en la oscuridad a partir de las reservas,
semanas antes de que la plántula se alimente por sí misma --- y, con la
cubierta, protección mientras tanto; los \qty{1.6e-4}{J} de la espora
construyen unas pocas células que deben hacer la fotosíntesis de
inmediato.
\textbf{24.} Semilla: fecundación sin agua; un embrión aprovisionado,
protegido y latente que se establece en lugares secos o estacionales;
dispersión por alas, frutos y animales. La estrategia del \omterm{def:b2:plant-life-cycles:fern}{helecho} gana
en hábitats húmedos, sombríos y estables, y en la colonización de
terrenos lejanos o recién abiertos, donde importan más los propágulos
baratos, ligeros y numerosos que el aprovisionamiento.
\textbf{25.} \omterm{def:b2:plant-life-cycles:fern}{Helecho}: $5.1\times 10^{6}$ esporas por temporada para 1,9
adultos, $2.7\times 10^{6}$ esporas por descendiente adulto, \qty{430}{J}
cada uno; pino: $9.6\times 10^{5}$ semillas para 480 adultos, \num{2000}
semillas por descendiente adulto, \qty{100}{kJ} cada uno.
\end{solution}
